% АВТОГЕНЕРОВАНО tools/build_index.py — не редагувати вручну
\multicolumn{6}{l}{\rule{0pt}{2.6ex}\bfseries\sffamily Батьківський реєстр встановленого (E) --- 40}\\*
\textbf{E1} & Відображення $\psi\to$ скаляри не ліпшицеве в $L^2$ & встановлено; уточнено ($\to$E36) & \nolinkurl{ESTABLISHED.md}\newline р.~23 & \nolinkurl{ourtry/03-deep-dives/D1-psi-to-scalars/conditioning.py:scalars_of} & \ref{ch:g02}\\
\textbf{E2} & Показники $\alpha$ розпадаються на три кластери & встановлено; лише 3 розряди ($\to$R10) & \nolinkurl{ESTABLISHED.md}\newline р.~24 & \nolinkurl{ourtry/03-deep-dives/D1-psi-to-scalars/conditioning.py:main} & \ref{ch:g02}\\
\textbf{E3} & Хвіст PCA шкідливіший за білий шум тієї ж енергії & встановлено (виміряно) & \nolinkurl{ESTABLISHED.md}\newline р.~25 & \nolinkurl{ourtry/03-deep-dives/D1-psi-to-scalars/conditioning.py:pert_pca_tail} & \ref{ch:g02}\\
\textbf{E4} & $L^2$-середнє двох гілок не є розв'язком & встановлено (виміряно) & \nolinkurl{ESTABLISHED.md}\newline р.~26 & \nolinkurl{ourtry/03-deep-dives/D2-gs-nonuniqueness/branch_mean.py:rel_residual_bratu} & \ref{ch:g03}\\
\textbf{E5} & Точка складки Братý $\lambda^*$ відтворена до останньої цифри & встановлено (виміряно) & \nolinkurl{ESTABLISHED.md}\newline р.~27 & \nolinkurl{ourtry/03-deep-dives/D2-gs-nonuniqueness/branch_mean.py:lam_star} & \ref{ch:g03}\\
\textbf{E6} & GS-нев'язка обчислювана лише з поданого $\psi$ & встановлено (виміряно) & \nolinkurl{ESTABLISHED.md}\newline р.~28 & \nolinkurl{ourtry/04-novelty/gs_residual_probe.py:gs_inconsistency} & \ref{ch:g04}\\
\textbf{E7} & GS-член інваріантний до масштабу $\psi$ & встановлено (виміряно) & \nolinkurl{ESTABLISHED.md}\newline р.~29 & \nolinkurl{ourtry/04-novelty/gs_residual_probe.py:gs_inconsistency} & \ref{ch:g04}\\
\textbf{E8} & GS-член сліпий до гладкої похибки й PCA-усічення (PCA-5: 0,0093) & встановлено; число уточнено & \nolinkurl{ESTABLISHED.md}\newline р.~30 & \nolinkurl{ourtry/04-novelty/pca5_check.py} & \ref{ch:g04}\\
\textbf{E9} & Скелет істини канонічний: 1 O-точка, 0 сідел & встановлено (виміряно) & \nolinkurl{ESTABLISHED.md}\newline р.~31 & \nolinkurl{ourtry/04-novelty/topology_probe.py:critical_points} & \ref{ch:g05}\\
\textbf{E10} & Скелет стійкий до 3\,\% шуму, руйнується на 10\,\% & встановлено; уточнено (N4) & \nolinkurl{ESTABLISHED.md}\newline р.~32 & \nolinkurl{ourtry/04-novelty/topology_probe.py:critical_points} & \ref{ch:g05}\\
\textbf{E11} & Сума індексів зберігається й у хаосі; інформативна лише кількість & встановлено; уточнено (N4) & \nolinkurl{ESTABLISHED.md}\newline р.~33 & \nolinkurl{ourtry/04-novelty/topology_probe.py:_winding} & \ref{ch:g05}\\
\textbf{E12} & На істині $\chi_{\mathrm{E}}(\psi\le t)\equiv1$ на всіх рівнях & встановлено (виміряно) & \nolinkurl{ESTABLISHED.md}\newline р.~34 & \nolinkurl{ourtry/04-novelty/topology_ec_curve.py:ec_curve} & \ref{ch:g05}\\
\textbf{E13} & Крива Ейлера менш чутлива за лічильник критичних точок & встановлено; уточнено (N4) & \nolinkurl{ESTABLISHED.md}\newline р.~35 & \nolinkurl{ourtry/04-novelty/topology_ec_curve.py:ec_curve} & \ref{ch:g05}\\
\textbf{E14} & Зміщення осі під шумом $\propto1/\sqrt{\lambda}$ (відтворено для $R$) & встановлено; відтворено 2026-09-30 & \nolinkurl{ESTABLISHED.md}\newline р.~36 & \nolinkurl{ourtry/03-deep-dives/D1-psi-to-scalars/c1/local_models.py:part_a} & \ref{ch:g02}\\
\textbf{E15} & Топологічний критерій машинно-незалежний (DIII-D і MAST) & встановлено (виміряно) & \nolinkurl{ESTABLISHED.md}\newline р.~37 & \nolinkurl{ourtry/04-novelty/h5_transfer_topology.py:skeleton} & \ref{ch:g06}\\
\textbf{E16} & Топологія переживає перенесення, значення $\psi$ --- ні & встановлено; примітка ($k=50$) & \nolinkurl{ESTABLISHED.md}\newline р.~38 & \nolinkurl{ourtry/04-novelty/h5_transfer_topology.py:skeleton} & \ref{ch:g06}\\
\textbf{E17} & Харнес скорера цілий: perfect $\to1$, zeros $\to0$ & встановлено (виміряно) & \nolinkurl{ESTABLISHED.md}\newline р.~39 & \nolinkurl{fec/starter/local_score.py:score_shot} & \emph{лише тут}\\
\textbf{E18} & В осесиметрії перетин Пуанкаре --- рівно контури $\psi$ & встановлено (прочитано) & \nolinkurl{ESTABLISHED.md}\newline р.~55 & \nolinkurl{viz/src/tokviz/fieldline.py:invariant_drift} & \ref{ch:g07}\\
\textbf{E19} & Імпульс --- тороїдальний потік, гамільтоніан --- полоїдальний & встановлено (прочитано) & \nolinkurl{ESTABLISHED.md}\newline р.~56 & \nolinkurl{viz/src/tokviz/perturbation.py:island_width_psin} & \ref{ch:g07}, \ref{ch:g08}\\
\textbf{E20} & Одна резонансна мода не дає хаосу & встановлено (прочитано) & \nolinkurl{ESTABLISHED.md}\newline р.~57 & --- & \ref{ch:g07}\\
\textbf{E21} & Розв'язок GS може бути неєдиним на реальній геометрії & встановлено (прочитано) & \nolinkurl{ESTABLISHED.md}\newline р.~58 & --- & \ref{ch:g03}\\
\textbf{E22} & Мала $L^2$-похибка не тягне малого фізичного резидуалу & встановлено (прочитано) & \nolinkurl{ESTABLISHED.md}\newline р.~59 & --- & \ref{ch:g04}\\
\textbf{E23} & Жоден фузійний бенчмарк не оцінює резидуал PDE & встановлено (прочитано) & \nolinkurl{ESTABLISHED.md}\newline р.~60 & --- & \ref{ch:g04}\\
\textbf{E24} & Стан Біна опуклий і однозначний; задача плазми --- ні & встановлено (прочитано) & \nolinkurl{ESTABLISHED.md}\newline р.~61 & --- & \ref{ch:g03}\\
\textbf{E25} & Прогноз зриву насичений і має закриті дані & встановлено (прочитано) & \nolinkurl{ESTABLISHED.md}\newline р.~62 & --- & \ref{ch:g01}\\
\textbf{E26} & Корпус Sophelio: 103,86~ГБ, 9\,121 розряд & встановлено (прочитано) & \nolinkurl{ESTABLISHED.md}\newline р.~63 & --- & \ref{ch:g01}\\
\textbf{E27} & Consistency усереднює сім скалярів, не вісім & встановлено (прочитано) & \nolinkurl{ESTABLISHED.md}\newline р.~64 & --- & \ref{ch:g01}\\
\textbf{E28} & Ліцензії HDB5, ConStellaration, Sophelio, FAIR-MAST & встановлено (прочитано) & \nolinkurl{ESTABLISHED.md}\newline р.~65 & --- & \ref{ch:g01}, \ref{ch:g06}\\
\textbf{E29} & Код starter kit справді MIT & встановлено (прочитано) & \nolinkurl{ESTABLISHED.md}\newline р.~66 & --- & \ref{ch:g01}\\
\textbf{E30} & Tokamap реалізовано; він точно симплектичний & встановлено (виміряно) & \nolinkurl{ESTABLISHED.md}\newline р.~40 & \nolinkurl{ourtry/03-deep-dives/D3-poincare-inverse/tokamap.py:jacobian_det} & \ref{ch:g07}\\
\textbf{E31} & Спостережуване на одній фазі непридатне; «540$\times$» порівнює різні стенди & встановлено; із застереженням & \nolinkurl{ESTABLISHED.md}\newline р.~41 & \nolinkurl{ourtry/03-deep-dives/D3-poincare-inverse/identifiability.py:profile} & \ref{ch:g07}\\
\textbf{E32} & Відновлення однієї амплітуди тривіальне (SNR 57) & встановлено (виміряно) & \nolinkurl{ESTABLISHED.md}\newline р.~42 & \nolinkurl{ourtry/03-deep-dives/D3-poincare-inverse/identifiability.py:main} & \ref{ch:g07}\\
\textbf{E33} & Функція Гріна GS-solver брала $k$ замість $m=k^2$ & виправлено & \nolinkurl{ESTABLISHED.md}\newline р.~43 & \nolinkurl{gs/src/numerical/compute.py:GreenFunction} & \ref{ch:g04}, \ref{ch:g10}\\
\textbf{E34} & GS-solver не запускався з README; полагоджено частково & виправлено частково & \nolinkurl{ESTABLISHED.md}\newline р.~44 & \nolinkurl{gs/src/numerical/grad_shafranov.py:GSsolverFreeBoundary} & \ref{ch:g10}\\
\textbf{E35} & Корінь R8: поправка $T'$ бере первісну $h(T)$, а не $V'(T)$ & встановлено (виміряно) & \nolinkurl{ESTABLISHED.md}\newline р.~45 & \nolinkurl{ourtry/03-deep-dives/D3-poincare-inverse/multimode_fixed_check.py:step_fixed} & \ref{ch:g07}\\
\textbf{E36} & Похибка осі: нахил 1 при $\eta<\eta^*=\lambda h^2/2$, $\approx\tfrac12$ вище & встановлено (виміряно) & \nolinkurl{ESTABLISHED.md}\newline р.~46 & \nolinkurl{ourtry/03-deep-dives/D1-psi-to-scalars/c1/local_models.py:part_a} & \ref{ch:g02}\\
\textbf{E37} & UNet: $R^2_\psi=0{,}976$, але Consistency 0,179 < PCA+Ridge 0,270 & встановлено (виміряно) & \nolinkurl{ESTABLISHED.md}\newline р.~47 & \nolinkurl{ourtry/03-deep-dives/D1-psi-to-scalars/c2/evaluate.py:main} & \ref{ch:g02}\\
\textbf{E38} & Мережі: $R^2_\psi\approx0{,}97$, але $g=0{,}81$--$0{,}96$ --- не рівноваги & встановлено (виміряно) & \nolinkurl{ESTABLISHED.md}\newline р.~48 & \nolinkurl{ourtry/04-novelty/c4/evaluate.py:main} & \ref{ch:g04}\\
\textbf{E39} & Вкладені поверхні без осесиметрії можливі (Landreman 2026; перевірено власним кодом) & встановлено (прочитано) & \nolinkurl{ESTABLISHED.md}\newline р.~67 & \nolinkurl{ourtry/99-bibliography/landreman2026/verify_landreman.py:check} & \ref{ch:g07}\\
\textbf{E40} & Ландшафт нев'язки екскурсії «скляний»: базовий фіт застрягає біля істини & встановлено (виміряно) & \nolinkurl{ESTABLISHED.md}\newline р.~49 & \nolinkurl{ourtry/03-deep-dives/D3-poincare-inverse/c5/stand_tokamap.py:observable} & \ref{ch:g07}\\
\multicolumn{6}{l}{\rule{0pt}{2.6ex}\bfseries\sffamily Спростоване батьківського реєстру (R) --- 13}\\*
\textbf{R1} & «Ніхто не нав'язує інтегральну крайову умову жорстко» & спростовано & \nolinkurl{ESTABLISHED.md}\newline р.~79 & --- & \ref{ch:g01}\\
\textbf{R2} & «Бін і вільна межа плазми --- одна математика» & спростовано & \nolinkurl{ESTABLISHED.md}\newline р.~80 & --- & \ref{ch:g03}\\
\textbf{R3} & «Симплектична нейромережа для силових ліній --- наша ідея» & спростовано & \nolinkurl{ESTABLISHED.md}\newline р.~81 & --- & \ref{ch:g07}\\
\textbf{R4} & «$\psi$ як канонічний імпульс»; повторилося в коді viz & спростовано; рецидив виправлено & \nolinkurl{ESTABLISHED.md}\newline р.~82 & \nolinkurl{viz/src/tokviz/perturbation.py:island_width_psin} & \ref{ch:g07}, \ref{ch:g08}\\
\textbf{R5} & «Крива Ейлера краща за лічильник критичних точок» & спростовано & \nolinkurl{ESTABLISHED.md}\newline р.~83 & \nolinkurl{ourtry/04-novelty/topology_ec_curve.py:ec_curve} & \ref{ch:g05}\\
\textbf{R6} & «$\alpha=\tfrac12$ --- універсальний закон» (механізм відтворено) & спростовано; механізм відтворено ($\to$E36) & \nolinkurl{ESTABLISHED.md}\newline р.~84 & \nolinkurl{ourtry/03-deep-dives/D1-psi-to-scalars/c1/local_models.py:part_a} & \ref{ch:g02}\\
\textbf{R7} & «Провал перенесення має топологічну складову» (H5) & спростовано & \nolinkurl{ESTABLISHED.md}\newline р.~85 & \nolinkurl{ourtry/04-novelty/h5_transfer_topology.py:skeleton} & \ref{ch:g06}\\
\textbf{R8} & «Tokamap переноситься на суму мод підстановкою» & спростовано; корінь знайдено ($\to$E35) & \nolinkurl{ESTABLISHED.md}\newline р.~86 & \nolinkurl{ourtry/03-deep-dives/D3-poincare-inverse/multimode.py:det_J} & \ref{ch:g07}\\
\textbf{R9} & «Провали LCFS MAST --- систематичний сигнал» (колишня C8) & спростовано (артефакт; $\leftarrow$C8) & \nolinkurl{ESTABLISHED.md}\newline р.~87 & \nolinkurl{ourtry/04-novelty/c8_sign_check.py} & \ref{ch:g06}\\
\textbf{R10} & «Три кластери $\alpha$ --- три механізми» (колишня C1) & спростовано ($\leftarrow$C1) & \nolinkurl{ESTABLISHED.md}\newline р.~88 & \nolinkurl{ourtry/03-deep-dives/D1-psi-to-scalars/c1/real_sweep.py:run_machine} & \ref{ch:g02}\\
\textbf{R11} & «Спектральна вага з $S(k)$ покращує моделі» (колишня C2) & спростовано ($\leftarrow$C2) & \nolinkurl{ESTABLISHED.md}\newline р.~89 & \nolinkurl{ourtry/03-deep-dives/D1-psi-to-scalars/c2/train.py:unet} & \ref{ch:g02}\\
\textbf{R12} & «GS\_score змінить упорядкування бейзлайнів» (половина C4) & спростовано ($\leftarrow$половина C4) & \nolinkurl{ESTABLISHED.md}\newline р.~90 & \nolinkurl{ourtry/04-novelty/c4/evaluate.py:main} & \ref{ch:g04}\\
\textbf{R13} & «Обернена задача Рівня~3 трактабельна за 48~год» (колишня C5) & спростовано & \nolinkurl{ESTABLISHED.md}\newline р.~91 & \nolinkurl{ourtry/03-deep-dives/D3-poincare-inverse/c5/fit_common.py:run_one} & \ref{ch:g07}, \ref{ch:g10}\\
\multicolumn{6}{l}{\rule{0pt}{2.6ex}\bfseries\sffamily Гіпотези (C) --- 10}\\*
\textbf{C1} & Три кластери $\alpha$ відповідають трьом механізмам ($\to$R10) & закрито: спростовано ($\to$R10) & \nolinkurl{CONJECTURES.md}\newline р.~21 & \nolinkurl{ourtry/03-deep-dives/D1-psi-to-scalars/c1/real_sweep.py:run_machine} & \ref{ch:g02}, \ref{ch:g10}\\
\textbf{C2} & Спектрально зважена втрата покращує реальні моделі ($\to$R11) & закрито: спростовано ($\to$R11) & \nolinkurl{CONJECTURES.md}\newline р.~22 & \nolinkurl{ourtry/03-deep-dives/D1-psi-to-scalars/c2/train.py:unet} & \ref{ch:g02}, \ref{ch:g10}\\
\textbf{C3} & Число обумовленості $1/\sqrt{\lambda}$ придатне для стратифікації тесту & активна; E14 відновлено, $\lambda$ широкий & \nolinkurl{CONJECTURES.md}\newline р.~23 & --- & \ref{ch:g02}, \ref{ch:g10}\\
\textbf{C4} & Калібрування змінить упорядкування бейзлайнів (GS-половина $\to$R12) & активна лише Calibration (GS $\to$R12) & \nolinkurl{CONJECTURES.md}\newline р.~24 & \nolinkurl{ourtry/04-novelty/c4/evaluate.py:main} & \ref{ch:g04}, \ref{ch:g10}\\
\textbf{C5} & Обернена задача Рівня~3 трактабельна за 48~год ($\to$R13) & закрито: спростовано ($\to$R13) & \nolinkurl{CONJECTURES.md}\newline р.~25 & \nolinkurl{viz/src/run_inverse.py:observable} & \ref{ch:g07}, \ref{ch:g10}\\
\textbf{C6} & Одночасні конформні смуги --- правильна форма UQ-члена & активна & \nolinkurl{CONJECTURES.md}\newline р.~26 & --- & \ref{ch:g04}, \ref{ch:g10}\\
\textbf{C7} & Deep Ritz нестабільний на плазмі, стабільний на Біні & активна & \nolinkurl{CONJECTURES.md}\newline р.~27 & --- & \ref{ch:g03}, \ref{ch:g10}\\
\textbf{C8} & Провали вилучення LCFS на MAST --- систематичний сигнал ($\to$R9) & закрито: спростовано ($\to$R9) & \nolinkurl{CONJECTURES.md}\newline р.~28 & \nolinkurl{ourtry/04-novelty/h5_transfer_topology.py:skeleton} & \ref{ch:g06}, \ref{ch:g10}\\
\textbf{C9} & Коваріантне подання покращить перенесення між машинами & відкладено & \nolinkurl{CONJECTURES.md}\newline р.~58 & --- & \ref{ch:g06}, \ref{ch:g10}\\
\textbf{C10} & Частина похибки моделей --- успадковане зміщення EFIT & відкладено & \nolinkurl{CONJECTURES.md}\newline р.~59 & --- & \ref{ch:g06}, \ref{ch:g10}\\
\multicolumn{6}{l}{\rule{0pt}{2.6ex}\bfseries\sffamily Відкриті питання (Q) --- 13}\\*
\textbf{Q1} & Чи реалістична реконструкція $\psi$ без магнітної діагностики & відкрите & \nolinkurl{open-questions.md}\newline р.~12 & --- & \ref{ch:g10}\\
\textbf{Q2} & Наскільки виродження $p'$/$FF'$ псує реконструкцію & відкрите & \nolinkurl{open-questions.md}\newline р.~13 & --- & \ref{ch:g10}\\
\textbf{Q3} & Чи трапляється неєдиність GS у робочих режимах & відкрите & \nolinkurl{open-questions.md}\newline р.~14 & --- & \ref{ch:g03}, \ref{ch:g10}\\
\textbf{Q4} & Реалістичні $\tau$ і $\Delta H$ для NbTi/Nb$_3$Sn & відкрите & \nolinkurl{open-questions.md}\newline р.~15 & --- & \ref{ch:g10}\\
\textbf{Q5} & Наскільки камера екранує швидке $\dd B/\dd t$ & відкрите & \nolinkurl{open-questions.md}\newline р.~16 & --- & \ref{ch:g10}\\
\textbf{Q6} & Доступ до кластера AI-лабораторії & відкрите & \nolinkurl{open-questions.md}\newline р.~22 & --- & \ref{ch:g10}\\
\textbf{Q7} & Дати, місце, розмір команди, журі & відкрите & \nolinkurl{open-questions.md}\newline р.~23 & --- & \ref{ch:g10}\\
\textbf{Q8} & Чи потрібна платформа зі скорингом & відкрите & \nolinkurl{open-questions.md}\newline р.~24 & --- & \ref{ch:g10}\\
\textbf{Q9} & Хто реально будує starter kit і скільки годин & відкрите & \nolinkurl{open-questions.md}\newline р.~25 & --- & \ref{ch:g10}\\
\textbf{Q10} & Умови поширення даних \texttt{fusionsimulator.io} & відкрите; не публікувати похідне & \nolinkurl{open-questions.md}\newline р.~31 & --- & \ref{ch:g01}, \ref{ch:g10}\\
\textbf{Q11} & Ліцензія ITPA HDB5 (CC BY 4.0) & закрите & \nolinkurl{open-questions.md}\newline р.~32 & --- & \ref{ch:g01}\\
\textbf{Q12} & Атрибуція скорера й даних Sophelio & закрите & \nolinkurl{open-questions.md}\newline р.~33 & --- & \ref{ch:g01}\\
\textbf{Q13} & Перевипуск MAST із CC BY-SA у CC BY --- підстава? & відкрите; блокує перевипуск & \nolinkurl{open-questions.md}\newline р.~34 & --- & \ref{ch:g01}, \ref{ch:g06}, \ref{ch:g10}\\
\multicolumn{6}{l}{\rule{0pt}{2.6ex}\bfseries\sffamily Рішення (D) --- 9}\\*
\textbf{D1} & Ядро --- реконструкція $\psi$, а не прогноз зриву & прийнято & \nolinkurl{DECISIONS.md}\newline р.~14 & --- & \ref{ch:g01}\\
\textbf{D2} & Наскрізний конвеєр --- лише в наративі & прийнято & \nolinkurl{DECISIONS.md}\newline р.~15 & --- & \ref{ch:g01}\\
\textbf{D3} & Два треки: базовий і дослідницький & прийнято & \nolinkurl{DECISIONS.md}\newline р.~16 & --- & \ref{ch:g01}\\
\textbf{D4} & Дані --- гібрид sim-to-real & прийнято & \nolinkurl{DECISIONS.md}\newline р.~17 & --- & \ref{ch:g01}\\
\textbf{D5} & Синтетика з FreeGS/FreeGSNKE/TokaLab, без власного генератора & прийнято & \nolinkurl{DECISIONS.md}\newline р.~18 & --- & \ref{ch:g01}\\
\textbf{D6} & Метрика $S'$: Sophelio + GS-резидуал + калібрування UQ & прийнято & \nolinkurl{DECISIONS.md}\newline р.~19 & --- & \ref{ch:g01}, \ref{ch:g04}\\
\textbf{D7} & Артефакти двомовні & прийнято & \nolinkurl{DECISIONS.md}\newline р.~20 & --- & \ref{ch:g01}\\
\textbf{D8} & Надпровідність --- лише як опуклий контроль (Бін) у D2 & прийнято; закрито & \nolinkurl{DECISIONS.md}\newline р.~21 & --- & \ref{ch:g01}, \ref{ch:g03}\\
\textbf{D9} & Квенч магніта --- не оцінювана задача & прийнято & \nolinkurl{DECISIONS.md}\newline р.~22 & --- & \ref{ch:g01}\\
\multicolumn{6}{l}{\rule{0pt}{2.6ex}\bfseries\sffamily Журнал верифікації (V-) --- 17}\\*
\textbf{V-A1} & Deep Ritz до рівняння GS не застосовували (пошук arXiv) & підтверджено (попередньо) & \nolinkurl{verification-log.md}\newline р.~20 & --- & \emph{лише тут}\\
\textbf{V-A2} & Жорстка інтегральна умова вже є (McClenaghan 2024; $\to$R1) & спростовано & \nolinkurl{verification-log.md}\newline р.~48 & --- & \emph{лише тут}\\
\textbf{V-A3} & Патологію «середнього гілок» для GS не квантифіковано & підтверджено (попередньо) & \nolinkurl{verification-log.md}\newline р.~68 & --- & \emph{лише тут}\\
\textbf{V-A4} & Жоден фузійний ML-бенчмарк не оцінює резидуал GS ($\to$E23) & підтверджено із застереженнями & \nolinkurl{verification-log.md}\newline р.~91 & --- & \emph{лише тут}\\
\textbf{V-D1} & Симплектична мережа для силових ліній --- HénonNet ($\to$R3) & спростовано & \nolinkurl{verification-log.md}\newline р.~121 & --- & \emph{лише тут}\\
\textbf{V-D2} & Топологічна вісь: новизна застосування, не методу & частково існує & \nolinkurl{verification-log.md}\newline р.~140 & --- & \emph{лише тут}\\
\textbf{V-D3} & Канонічні змінні: «$\psi$ як імпульс» хибне ($\to$R4) & чернетку спростовано & \nolinkurl{verification-log.md}\newline р.~163 & --- & \ref{ch:g08}\\
\textbf{V-D4} & Власні виміри топології ($\to$E9--E13, R5) & перевірено (власні виміри) & \nolinkurl{verification-log.md}\newline р.~176 & --- & \emph{лише тут}\\
\textbf{V-B1} & Обсяг корпусу: 103,86~ГБ ($\to$E26) & перевірено & \nolinkurl{verification-log.md}\newline р.~194 & --- & \emph{лише тут}\\
\textbf{V-B2} & Кількість розрядів: 9\,121 ($\to$E26) & перевірено & \nolinkurl{verification-log.md}\newline р.~195 & --- & \emph{лише тут}\\
\textbf{V-B3} & Consistency: сім скалярів ($\to$E27) & перевірено & \nolinkurl{verification-log.md}\newline р.~196 & --- & \emph{лише тут}\\
\textbf{V-B4} & Канонічний репозиторій starter kit --- Sophelio & перевірено & \nolinkurl{verification-log.md}\newline р.~197 & --- & \emph{лише тут}\\
\textbf{V-B5} & Топові бали недоступні без облікового запису & перевірено & \nolinkurl{verification-log.md}\newline р.~198 & --- & \emph{лише тут}\\
\textbf{V-C1} & ITPA HDB5 --- CC BY 4.0 ($\to$Q11) & перевірено [V] & \nolinkurl{verification-log.md}\newline р.~206 & --- & \emph{лише тут}\\
\textbf{V-C2} & ConStellaration --- MIT & перевірено [V] & \nolinkurl{verification-log.md}\newline р.~207 & --- & \emph{лише тут}\\
\textbf{V-C3} & Sophelio FEC --- CC BY 4.0 & перевірено [V] & \nolinkurl{verification-log.md}\newline р.~208 & --- & \emph{лише тут}\\
\textbf{V-C4} & FAIR-MAST --- CC BY-SA 4.0 ($\to$Q13) & перевірено [V] & \nolinkurl{verification-log.md}\newline р.~209 & --- & \emph{лише тут}\\
\multicolumn{6}{l}{\rule{0pt}{2.6ex}\bfseries\sffamily Візуалізація: встановлене (VE) --- 34}\\*
\textbf{VE1} & Пошук осі збігається з EFIT до 0,0002~мм & встановлено (виміряно) & \nolinkurl{findings.md}\newline р.~26 & \nolinkurl{viz/src/tokviz/equilibrium.py:find_axis} & \ref{ch:g08}\\
\textbf{VE2} & LCFS набору даних --- справді контур $\psi$ & встановлено (виміряно) & \nolinkurl{findings.md}\newline р.~31 & \nolinkurl{viz/src/tokviz/equilibrium.py:set_boundary_from_contour} & \ref{ch:g08}\\
\textbf{VE3} & $F$ з $q_{95}$ дає $B_\varphi=1{,}927$~Тл без знання $B_0$ & встановлено (виміряно) & \nolinkurl{findings.md}\newline р.~35 & \nolinkurl{viz/src/tokviz/equilibrium.py:calibrate_F} & \ref{ch:g08}\\
\textbf{VE4} & $q$ контурним інтегралом і трасуванням збігаються до 0,16\,\% & встановлено (виміряно) & \nolinkurl{findings.md}\newline р.~41 & \nolinkurl{viz/src/tokviz/fieldline.py:q_from_tracing} & \ref{ch:g08}\\
\textbf{VE5} & Інтегратор зберігає $\psi$ до $10^{-7}$ за 200 обертів & встановлено (виміряно) & \nolinkurl{findings.md}\newline р.~46 & \nolinkurl{viz/src/tokviz/fieldline.py:invariant_drift} & \ref{ch:g08}\\
\textbf{VE6} & Растеризація $\cos m\theta^*$ на 65$\times$65 дає фальшиві острови & встановлено (виміряно) & \nolinkurl{findings.md}\newline р.~51 & \nolinkurl{viz/src/tokviz/perturbation.py:Perturbation} & \ref{ch:g08}\\
\textbf{VE7} & Аналітичні похідні збурення збігаються зі скінченними різницями & встановлено (виміряно) & \nolinkurl{findings.md}\newline р.~60 & \nolinkurl{viz/tests/test_physics.py:test_perturbation_derivatives_match_finite_difference} & \ref{ch:g08}\\
\textbf{VE8} & Ланцюг $m/n$ дає рівно $m$ островів у розрізі & встановлено (виміряно) & \nolinkurl{findings.md}\newline р.~63 & \nolinkurl{viz/src/tokviz/islands.py:classify_fixed_points} & \ref{ch:g08}\\
\textbf{VE9} & Чириков $S=0{,}458<1$ (було 0,714); поверхні виживають & встановлено (виміряно) & \nolinkurl{findings.md}\newline р.~69 & \nolinkurl{viz/src/tokviz/perturbation.py:chirikov} & \ref{ch:g08}\\
\textbf{VE10} & Перший рендер Cycles на Metal: $\sim$100~с компіляції & встановлено (виміряно) & \nolinkurl{findings.md}\newline р.~75 & --- & \emph{лише тут}\\
\textbf{VE11} & Продуктивність Cycles на M4: GPU лише $\sim$1,2$\times$ & встановлено (виміряно) & \nolinkurl{findings.md}\newline р.~86 & --- & \emph{лише тут}\\
\textbf{VE11a} & Час рендера з тестової сцени занижено вдвічі & встановлено (виміряно) & \nolinkurl{findings.md}\newline р.~79 & --- & \emph{лише тут}\\
\textbf{VE12} & Blender не має імпортера STEP/IGES & встановлено (прочитано) & \nolinkurl{findings.md}\newline р.~250 & --- & \emph{лише тут}\\
\textbf{VE13} & Зміни API Blender 4.2$\to$5.x, що ламають скрипти & встановлено (прочитано) & \nolinkurl{findings.md}\newline р.~252 & --- & \emph{лише тут}\\
\textbf{VE14} & EEVEE: об'ємна емісія не освітлює поверхні & встановлено (прочитано) & \nolinkurl{findings.md}\newline р.~257 & --- & \emph{лише тут}\\
\textbf{VE15} & ITER-2024: перша стінка --- вольфрам замість берилію & встановлено (прочитано) & \nolinkurl{findings.md}\newline р.~260 & --- & \emph{лише тут}\\
\textbf{VE16} & Число обертання трасуванням відтворює $1/q$ до 0,03\,\% & встановлено (виміряно) & \nolinkurl{findings.md}\newline р.~90 & \nolinkurl{viz/src/tokviz/analysis.py:rotation_number} & \ref{ch:g08}\\
\textbf{VE17} & Хаотичні орбіти групуються рівно на резонансах & встановлено (виміряно) & \nolinkurl{findings.md}\newline р.~95 & \nolinkurl{viz/src/tokviz/analysis.py:classify} & \ref{ch:g08}\\
\textbf{VE18} & Виправлена маятникова ширина збігається з виміряною до 0--5\,\% & встановлено (виміряно) & \nolinkurl{findings.md}\newline р.~100 & \nolinkurl{viz/src/tokviz/analysis.py:measured_island_width} & \ref{ch:g08}\\
\textbf{VE19} & Хаос існує задовго до $S=1$; перехід через 1 не перевірено & встановлено (виміряно) & \nolinkurl{findings.md}\newline р.~136 & \nolinkurl{viz/src/tokviz/analysis.py:classify} & \ref{ch:g08}\\
\textbf{VE20} & Згущення засіву завищує частку хаосу в 1,6--1,9 раза & встановлено (виміряно) & \nolinkurl{findings.md}\newline р.~130 & \nolinkurl{viz/src/run_analysis.py:main} & \ref{ch:g08}\\
\textbf{VE21} & $\Bvec=\nabla\times(\psi\nabla\varphi)$ бездивергентне для будь-якого числа мод & встановлено (виміряно) & \nolinkurl{findings.md}\newline р.~151 & \nolinkurl{viz/src/tokviz/fieldline.py:FieldLines} & \ref{ch:g07}\\
\textbf{VE22} & Відновлення самих амплітуд тривіальне (обумовленість 1,2) & встановлено (виміряно) & \nolinkurl{findings.md}\newline р.~159 & \nolinkurl{viz/src/run_inverse.py:observable} & \ref{ch:g07}\\
\textbf{VE23} & Складність оберненої задачі задає спостережуване, а не фізика & встановлено (виміряно) & \nolinkurl{findings.md}\newline р.~164 & \nolinkurl{viz/src/run_inverse.py:observable} & \ref{ch:g07}\\
\textbf{VE24} & Перевага роздільності = фазова інформація + обсяг даних & встановлено (виміряно) & \nolinkurl{findings.md}\newline р.~174 & \nolinkurl{viz/src/run_inverse.py:main} & \ref{ch:g07}\\
\textbf{VE25} & Затиснута мода 5/2 гірша лише без фазової інформації & встановлено (виміряно) & \nolinkurl{findings.md}\newline р.~181 & \nolinkurl{viz/src/run_inverse.py:main} & \ref{ch:g07}\\
\textbf{VE26} & Паразитні гармоніки стенда не дають видимих островів & встановлено (виміряно) & \nolinkurl{findings.md}\newline р.~187 & \nolinkurl{viz/src/run_parasitic.py:main} & \ref{ch:g09}\\
\textbf{VE27} & Біо--Савар перевірено проти аналітичної петлі & встановлено (виміряно) & \nolinkurl{findings.md}\newline р.~197 & \nolinkurl{viz/src/tokviz/rmp_coils.py:biot_savart} & \ref{ch:g09}\\
\textbf{VE28} & Поле масиву $4\cdot10^{-23}$~Тл --- через точку симетрії & встановлено (виміряно) & \nolinkurl{findings.md}\newline р.~202 & \nolinkurl{viz/src/tokviz/rmp_coils.py:icoil_array} & \ref{ch:g09}\\
\textbf{VE29} & Інстанси Geometry Nodes не успадковують матеріали & встановлено (виміряно) & \nolinkurl{findings.md}\newline р.~207 & --- & \emph{лише тут}\\
\textbf{VE30} & Слід на стінці зважений до краю (VC6) & встановлено (виміряно) & \nolinkurl{findings.md}\newline р.~212 & \nolinkurl{viz/src/run_footprint.py:observable} & \ref{ch:g07}\\
\textbf{VE31} & Обумовленість гіршає з числом мод плавно (VC7) & встановлено (виміряно) & \nolinkurl{findings.md}\newline р.~220 & \nolinkurl{viz/src/run_mode_scan.py:main} & \ref{ch:g07}\\
\textbf{VE32} & Гістограма ударів непридатна для скінченних різниць & встановлено (виміряно) & \nolinkurl{findings.md}\newline р.~231 & \nolinkurl{viz/src/tokviz/footprint.py:connection_length_profile} & \emph{лише тут}\\
\textbf{VE33} & Немонотонність частки хаосу при 40 зернах --- шум (VC5) & встановлено (виміряно) & \nolinkurl{findings.md}\newline р.~236 & \nolinkurl{viz/src/run_analysis.py:main} & \ref{ch:g08}\\
\multicolumn{6}{l}{\rule{0pt}{2.6ex}\bfseries\sffamily Візуалізація: спростоване (VR) --- 16}\\*
\textbf{VR1} & «Геометричного кута досить для розміщення островів» & спростовано & \nolinkurl{findings.md}\newline р.~267 & \nolinkurl{viz/src/tokviz/equilibrium.py:theta_star_field} & \ref{ch:g08}\\
\textbf{VR2} & «Паразитні гармоніки інтерполяції можна ігнорувати» & спростовано & \nolinkurl{findings.md}\newline р.~271 & \nolinkurl{viz/src/tokviz/perturbation.py:Perturbation} & \ref{ch:g08}\\
\textbf{VR3} & «Об'ємну сітку плазми можна різати разом з машиною» & спростовано & \nolinkurl{findings.md}\newline р.~275 & --- & \emph{лише тут}\\
\textbf{VR4} & «Камери можна кадрувати за $R_0$» & спростовано & \nolinkurl{findings.md}\newline р.~279 & --- & \emph{лише тут}\\
\textbf{VR5} & «Спільна ціль TRACK\_TO для всіх камер» & спростовано & \nolinkurl{findings.md}\newline р.~283 & --- & \emph{лише тут}\\
\textbf{VR6} & «Можу навести замкнену форму кривої Princeton D» & спростовано & \nolinkurl{findings.md}\newline р.~286 & \nolinkurl{viz/src/tokviz/surfaces.py:d_shape} & \emph{лише тут}\\
\textbf{VR7} & «\texttt{angle2 == 0} означає зсув $0^\circ$» & спростовано & \nolinkurl{findings.md}\newline р.~291 & --- & \emph{лише тут}\\
\textbf{VR8} & «Число обертання --- просте середнє за вікном» & спростовано & \nolinkurl{findings.md}\newline р.~295 & \nolinkurl{viz/src/tokviz/analysis.py:rotation_number} & \ref{ch:g08}\\
\textbf{VR9} & «Острів розпізнається за екскурсією $\psi_N$» & спростовано & \nolinkurl{findings.md}\newline р.~301 & \nolinkurl{viz/src/tokviz/analysis.py:classify} & \ref{ch:g08}\\
\textbf{VR10} & «Захоплення $\nu$ --- близькість до раціонального» & спростовано & \nolinkurl{findings.md}\newline р.~307 & \nolinkurl{viz/src/tokviz/analysis.py:find_locked_bands} & \ref{ch:g08}\\
\textbf{VR11} & «Класифікувати можна в будь-якому порядку» & спростовано & \nolinkurl{findings.md}\newline р.~312 & \nolinkurl{viz/src/tokviz/analysis.py:classify} & \ref{ch:g08}\\
\textbf{VR12} & «Дефіцит ширини острова --- дискретність засіву» & спростовано & \nolinkurl{findings.md}\newline р.~317 & \nolinkurl{viz/src/tokviz/analysis.py:measured_island_width} & \ref{ch:g08}\\
\textbf{VR13} & «Взаємної узгодженості $\nu$ досить для захоплення» & спростовано & \nolinkurl{findings.md}\newline р.~338 & \nolinkurl{viz/src/tokviz/analysis.py:find_locked_bands} & \ref{ch:g08}\\
\textbf{VR14} & Обвідна $\psi_N^{m/2}$ не відтворює поле реальних котушок & спростовано & \nolinkurl{findings.md}\newline р.~327 & \nolinkurl{viz/src/run_coil_compare.py:fit_exponent} & \ref{ch:g09}\\
\textbf{VR15} & Ширина острова з $\psi$ як імпульсом: завищення в $\sqrt q$ & спростовано; код виправлено & \nolinkurl{findings.md}\newline р.~345 & \nolinkurl{viz/src/tokviz/perturbation.py:island_width_psin} & \ref{ch:g08}\\
\textbf{VR16} & «Дефіцит ширини --- стохастичний шар сепаратриси» & спростовано & \nolinkurl{findings.md}\newline р.~361 & \nolinkurl{viz/tests/test_physics.py:test_island_width_matches_traced_separatrix} & \ref{ch:g08}\\
\multicolumn{6}{l}{\rule{0pt}{2.6ex}\bfseries\sffamily Візуалізація: гіпотези (VC) --- 7}\\*
\textbf{VC1} & Обвідна $\psi_N^{m/2}$ близька до поля RMP-котушок ($\to$VR14) & закрито: спростовано & \nolinkurl{findings.md}\newline р.~385 & \nolinkurl{viz/src/run_coil_compare.py:resonant_harmonic} & \ref{ch:g09}\\
\textbf{VC2} & Паразитні 1,3\,\% не дають видимих островів ($\to$VE26) & закрито: підтверджено & \nolinkurl{findings.md}\newline р.~392 & \nolinkurl{viz/src/run_parasitic.py:main} & \ref{ch:g09}\\
\textbf{VC3} & Шість обраних ракурсів достатні для лекції & відкрито & \nolinkurl{findings.md}\newline р.~410 & --- & \ref{ch:g10}\\
\textbf{VC4} & Кадри читатимуться при проєкції у великій залі & відкрито & \nolinkurl{findings.md}\newline р.~413 & --- & \ref{ch:g10}\\
\textbf{VC5} & Немонотонність частки хаосу зникає з числом зерен ($\to$VE33) & закрито: підтверджено & \nolinkurl{findings.md}\newline р.~401 & \nolinkurl{viz/src/run_analysis.py:main} & \ref{ch:g08}\\
\textbf{VC6} & Обернена задача на сліді важча за екскурсійну ($\to$VE30) & закрито: підтверджено & \nolinkurl{findings.md}\newline р.~372 & \nolinkurl{viz/src/run_footprint.py:main} & \ref{ch:g07}\\
\textbf{VC7} & Понад три моди обумовленість помітно гіршає ($\to$VE31) & закрито: підтверджено & \nolinkurl{findings.md}\newline р.~373 & \nolinkurl{viz/src/run_mode_scan.py:main} & \ref{ch:g07}\\
\multicolumn{6}{l}{\rule{0pt}{2.6ex}\bfseries\sffamily Візуалізація: відкрите (VQ) --- 5}\\*
\textbf{VQ1} & Кріостат моделі тороїдальний; в ITER --- циліндр із куполом & відкрите & \nolinkurl{findings.md}\newline р.~420 & --- & \ref{ch:g10}\\
\textbf{VQ2} & Конфлікт CC BY і CC BY-SA при перевипуску ($\leftarrow$Q13) & відкрите & \nolinkurl{findings.md}\newline р.~423 & --- & \ref{ch:g10}\\
\textbf{VQ3} & Чи додати виміряний перетин W7-X поруч із синтетичним & відкрите & \nolinkurl{findings.md}\newline р.~426 & --- & \ref{ch:g10}\\
\textbf{VQ4} & Чи надійний тест пласкості $\nu$ при новому засіві & відкрите & \nolinkurl{findings.md}\newline р.~429 & \nolinkurl{viz/src/tokviz/analysis.py:find_locked_bands} & \ref{ch:g08}, \ref{ch:g10}\\
\textbf{VQ5} & U1: 21/30 хаотичних проти 10--13/30 у таблиці --- неузгоджено & відкрите & \nolinkurl{findings.md}\newline р.~443 & \nolinkurl{viz/src/run_coil_compare.py} & \ref{ch:g09}\\
\multicolumn{6}{l}{\rule{0pt}{2.6ex}\bfseries\sffamily Бекенд котушок, пункти без номерів (U) --- 3}\\*
\textbf{U1} & «Хаос бекенда котушок спричиняє гребінка бічних смуг» & спростовано & \nolinkurl{COIL-BACKEND.md}\newline р.~91 & \nolinkurl{viz/src/tokviz/coilfield.py:CoilPerturbation} & \ref{ch:g09}\\
\textbf{U2} & Хаос бекенда котушок: фізика чи похибка реконструкції поля? & відкрите & \nolinkurl{COIL-BACKEND.md}\newline р.~95 & \nolinkurl{viz/src/tokviz/coilfield.py:CoilPerturbation} & \ref{ch:g09}, \ref{ch:g10}\\
\textbf{U3} & Тримати $\vect A$, а не $\Bvec$: $\nabla\cdot\Bvec$ з $3{,}19\cdot10^{-2}$ до $2{,}50\cdot10^{-6}$ & встановлено (виміряно) & \nolinkurl{COIL-BACKEND.md}\newline р.~35 & \nolinkurl{viz/src/tokviz/coilfield.py:divergence_report} & \ref{ch:g09}\\
